\documentclass[10pt,a4paper]{amsart}
\usepackage[margin=19mm]{geometry}
\usepackage{amsmath,amssymb,mathtools}
\usepackage[scr=rsfs,cal=euler]{mathalfa}
\usepackage{xfrac}
\usepackage[unicode,hidelinks,bookmarks=false]{hyperref}
\newcommand{\ie}{i.e.\ }
\newcommand{\cf}{cf.\ }
\newcommand{\resp}{respectively\ }
\newcommand{\Subsection}{\subsection}
\providecommand{\nobreakdash}{\nobreak\hbox{-}}
\setlength{\emergencystretch}{2em}
\multlinegap0pt
\usepackage{bm}
\usepackage{bbm}
\usepackage{dsfont}
\usepackage{xspace}
\usepackage{cancel}
\usepackage{mathtools}
\usepackage{amsthm}
\makeatletter
\@ifundefined{proposition}{\newtheorem{proposition}{Proposition}}{}
\makeatother
\newcommand{\tr}{\text{tr}}
\renewcommand{\div}{\mbox{div }}
\title[A large data result for vacuum Einstein's equations]{A large data result for vacuum Einstein's equations}
\author{Puskar Mondal}
\date{Source: arXiv:2502.11289v4 (CC BY 4.0)}
\begin{document}
\maketitle

\noindent\emph{Source and licence.} A large data result for vacuum Einstein's equations. Version arXiv:2502.11289v4 (CC BY 4.0), distributed under the Creative Commons Attribution 4.0 International licence, as recorded in the arXiv licence metadata for this exact version and shown on the abstract page https://arxiv.org/abs/2502.11289v4. Full rights notice: licenses/arxiv-2502.11289-CC-BY-4.0.txt.\par\smallskip

\noindent\textit{Editorial scope.} Counted unit: the Proposition whose source label is \texttt{prop:conformal-correction} in the article (source0.tex lines 3604--3657, source label \texttt{prop:conformal-correction}), delivered together with the article's own complete original proof of it (source0.tex lines 3659--4001, one \texttt{proof} environment of 343 lines). No mathematics is written, repaired or re-derived: statement and proof are the article's own text line for line. Required context retained uncounted: none. Every definition, display and label of those lines is retained unchanged. Omitted from this package and not redistributed: 1--3603, 3658--3658, 4002--5871. Nothing inside a retained excerpt is omitted or altered: every retained body is delivered exactly as the article's lines are written, so no omission receipt of any kind accompanies this package. Citations: the retained excerpts carry no citation command at all, so no bibliography entry of the article is transcribed and no reference is redistributed. Reference adaptation: none was needed and none was made; every reference inside the delivered unit resolves inside it, so no unresolved reference or citation is produced. Printed slips kept literal: no printed word, symbol, hypothesis or number of the counted statement or of its proof was altered. Layout and class: the article's own class is \texttt{amsart} with packages including pgfplots, amsmath, sansmath, graphicx, slashed, enumerate, dsfont, mathtools, cancel, footmisc, mathalfa, tikz-cd, hyperref, amsfonts; the wrapper is the lane's pinned \texttt{amsart} editorial wrapper at 10pt on A4, which restates the theorem-like environments the retained text needs and adds the article's own macro definitions that the retained text needs (\texttt{\textbackslash div}, \texttt{\textbackslash tr}), with extra wrapper packages (bm, bbm, dsfont, xspace, cancel, mathtools). The delivered numbering is the unit's own. Assets: no retained span contains a figure environment, a graphics-inclusion command, a PDF-page inclusion, a table or a TikZ picture; no image file is copied into this package, the package declares no asset dependency and no image file is redistributed with it. Licence: version arXiv:2502.11289v4 of this article is distributed under the Creative Commons Attribution 4.0 International licence, as recorded in the arXiv licence metadata for this exact version and shown on the abstract page https://arxiv.org/abs/2502.11289v4. A full rights notice is delivered at licenses/arxiv-2502.11289-CC-BY-4.0.txt. Characters: every retained excerpt is pure ASCII, so the delivered engine, XeLaTeX with its default Latin Modern text font, reports no missing character.
\begin{proposition}[Existence of a unique positive solution to the Lichnerowicz equation]
\label{prop:conformal-correction}
Fix a bounded-geometry class \(\mathcal G\) of smooth Riemannian metrics on the closed
\(3\)-manifold \(M\). Then there exist constants
\[
\delta_{0}>0,\qquad C\ge 1,
\]
depending only on \(\mathcal G\), with the following property.

\noindent Let \(g\in \mathcal G\), and let \(\Sigma\) be a smooth symmetric \(2\)-tensor satisfying
\[
\tr_{g}\Sigma=0,\qquad \div_{g}\Sigma=0,
\]
together with
\[
\left\|R[g]+\frac23\right\|_{H^{2}(g)}\le \delta_{0},
\qquad
\|\Sigma\|_{H^{2}(g)}\le \delta_{0}.
\]
Define
\[
\mathcal E(g,\Sigma):=R[g]+\frac23-|\Sigma|_{g}^{2}.
\]
Then there exists a unique positive solution
\[
\varphi\in H^{4}(M)
\]
of the Lichnerowicz equation
\begin{equation}
\label{eq:Lichnerowicz-exact-corrected}
-8\Delta_{g}\varphi + R[g]\varphi + \frac23 \varphi^{5} - |\Sigma|_{g}^{2}\varphi^{-7}=0.
\end{equation}
Moreover, \(\varphi\) is smooth, and if one defines
\[
g_{0}:=\varphi^{4}g,\qquad \Sigma_{0}:=\varphi^{-2}\Sigma,
\]
then
\begin{align}
R[g_{0}]+\frac23&=|\Sigma_{0}|_{g_{0}}^{2},\\
\div_{g_{0}}\Sigma_{0}&=0,\\
\tr_{g_{0}}\Sigma_{0}&=0.
\end{align}
In addition,
\begin{align}
\label{eq:phi-close-corrected}
\|\varphi-1\|_{H^{4}(g)}
&\leq
C\|\mathcal E(g,\Sigma)\|_{H^{2}(g)},\\
\label{eq:T-change-corrected}
\|\mathfrak T[g_{0}]-\mathfrak T[g]\|_{H^{2}(g)}
&\leq
C\|\mathcal E(g,\Sigma)\|_{H^{2}(g)}.
\end{align}
\end{proposition}

\begin{proof}
Set
\[
S:=|\Sigma|_{g}^{2},
\qquad
\rho:=R[g]+\frac23.
\]
Then define the following
\[
\mathcal E(g,\Sigma)=\rho-S
\]
which is small in $H^{2}$ norm. 
Since \(M\) is \(3\)-dimensional and \(g\in\mathcal G\), the Sobolev embeddings and
multiplication estimates are uniform in \(g\). In particular,
\[
\|S\|_{H^{2}(g)}\le C\|\Sigma\|_{H^{2}(g)}^{2},
\qquad
\|S\|_{C^{0}(g)}\le C\|\Sigma\|_{H^{2}(g)}^{2},
\qquad
\|\rho\|_{C^{0}(g)}\le C\|\rho\|_{H^{2}(g)}.
\]
Hence, for \(\delta_{0}\) sufficiently small,
\begin{equation}
\label{eq:defect-small-proof}
\|\mathcal E(g,\Sigma)\|_{H^{2}(g)}
\le C\delta_{0},
\qquad
\|\rho\|_{C^{0}(g)}+\|S\|_{C^{0}(g)}\le C\delta_{0}.
\end{equation}

\noindent We write
\[
\varphi=1+u
\]
and define
\[
\mathcal F_{g,\Sigma}(\varphi)
:=
-8\Delta_{g}\varphi + R[g]\varphi + \frac23\varphi^{5}-S\varphi^{-7}.
\]
Then \eqref{eq:Lichnerowicz-exact-corrected} is equivalent to
\[
\mathcal F_{g,\Sigma}(1+u)=0.
\]
Moreover,
\[
\mathcal F_{g,\Sigma}(1)=\rho-S=\mathcal E(g,\Sigma).
\]

\noindent First, we understand the linearization of $\mathcal F_{g,\Sigma}$. The linearization at \(u=0\) is
\[
L_{g,\Sigma}u
:=
D\mathcal F_{g,\Sigma}|_{\varphi=1}[u]
=
-8\Delta_{g}u+\Bigl(R[g]+\frac{10}{3}+7S\Bigr)u.
\]
Since
\[
R[g]+\frac{10}{3}+7S
=
\frac83+\rho+7S,
\]
the pointwise bound \eqref{eq:defect-small-proof} implies, after shrinking \(\delta_{0}\),
that
\[
R[g]+\frac{10}{3}+7S\ge 1
\qquad\text{on }M.
\]
Therefore
\[
\int_{M} \langle L_{g,\Sigma}u,u\rangle\,d\mu_{g}
=
8\|\nabla u\|_{L^{2}(g)}^{2}
+
\int_{M}\Bigl(R[g]+\frac{10}{3}+7S\Bigr)u^{2}\,d\mu_{g}
\ge c\|u\|_{H^{1}(g)}^{2}.
\]
Hence \(L_{g,\Sigma}\) has trivial kernel. Since it is a strongly elliptic second-order
operator on a closed manifold, it is Fredholm of index zero, and thus
\[
L_{g,\Sigma}\colon H^{4}(M)\to H^{2}(M)
\]
is an isomorphism. By uniform elliptic regularity on the bounded-geometry class
\(\mathcal G\), there is a constant \(K\) such that
\begin{equation}
\label{eq:elliptic-L-corrected}
\|u\|_{H^{4}(g)}
\le
K\|L_{g,\Sigma}u\|_{H^{2}(g)}.
\end{equation}

\noindent Next we isolate the nonlinear remainder. For \(u\) sufficiently small in \(C^{0}\),
\[
(1+u)^{5}=1+5u+Q_{5}(u),
\qquad
(1+u)^{-7}=1-7u+Q_{-7}(u),
\]
where \(Q_{5}(u)\) and \(Q_{-7}(u)\) vanish quadratically at \(u=0\). Thus
\[
\mathcal F_{g,\Sigma}(1+u)
=
\mathcal E(g,\Sigma)+L_{g,\Sigma}u+\mathcal N_{g,\Sigma}(u),
\]
with
\[
\mathcal N_{g,\Sigma}(u)
=
\frac23 Q_{5}(u)-S\,Q_{-7}(u).
\]

\noindent Let \(C_{\mathrm{emb}}\) denote the uniform norm of the embedding
\(H^{4}(M)\hookrightarrow C^{0}(M)\) on \(\mathcal G\), and set
\[
\rho_{*}:=\frac{1}{4C_{\mathrm{emb}}}.
\]
If \(\|u\|_{H^{4}(g)}\le \rho_{*}\), then \(\|u\|_{C^{0}(g)}\le \frac14\), so
\[
\frac34\le 1+u\le \frac54.
\]
On this set, uniform Moser estimates imply
\begin{align}
\label{eq:Nquad-corrected}
\|\mathcal N_{g,\Sigma}(u)\|_{H^{2}(g)}
&\le
C\|u\|_{H^{4}(g)}^{2},\\
\label{eq:Nlip-corrected}
\|\mathcal N_{g,\Sigma}(u)-\mathcal N_{g,\Sigma}(v)\|_{H^{2}(g)}
&\le
C\bigl(\|u\|_{H^{4}(g)}+\|v\|_{H^{4}(g)}\bigr)\|u-v\|_{H^{4}(g)}.
\end{align}

\noindent Define
\[
\Phi(u):=
-L_{g,\Sigma}^{-1}\bigl(\mathcal E(g,\Sigma)+\mathcal N_{g,\Sigma}(u)\bigr).
\]
Let
\[
r:=2K\|\mathcal E(g,\Sigma)\|_{H^{2}(g)}.
\]
If \(\delta_{0}\) is sufficiently small, then by \eqref{eq:defect-small-proof},
\[
r\le \rho_{*},
\qquad
CKr\le \frac12,
\qquad
CKr^{2}\le \frac r2.
\]
Now let \(\|u\|_{H^{4}(g)}\le r\). Using \eqref{eq:elliptic-L-corrected} and
\eqref{eq:Nquad-corrected},
\[
\|\Phi(u)\|_{H^{4}(g)}
\le
K\|\mathcal E(g,\Sigma)\|_{H^{2}(g)}
+
K\|\mathcal N_{g,\Sigma}(u)\|_{H^{2}(g)}
\le
\frac r2+CKr^{2}
\le r.
\]
Hence \(\Phi\) maps the closed \(H^{4}\)-ball \(B_{r}\) into itself. Likewise, for
\(u,v\in B_{r}\), \eqref{eq:Nlip-corrected} gives
\[
\|\Phi(u)-\Phi(v)\|_{H^{4}(g)}
\le
K\|\mathcal N_{g,\Sigma}(u)-\mathcal N_{g,\Sigma}(v)\|_{H^{2}(g)}
\le
2CKr\,\|u-v\|_{H^{4}(g)}
\le
\frac12\|u-v\|_{H^{4}(g)}.
\]
Thus \(\Phi\) is a contraction on \(B_{r}\). By Banach's fixed-point theorem, there exists
a unique \(u\in B_{r}\) such that \(\Phi(u)=u\). Setting \(\varphi:=1+u\), we obtain a
solution of \eqref{eq:Lichnerowicz-exact-corrected} with
\[
\|\varphi-1\|_{H^{4}(g)}=\|u\|_{H^{4}(g)}
\le
C\|\mathcal E(g,\Sigma)\|_{H^{2}(g)},
\]
which proves \eqref{eq:phi-close-corrected}. Since \(\|u\|_{C^{0}(g)}\le \frac14\), we also
have \(\varphi\ge \frac34>0\).

\noindent We now prove uniqueness among \emph{all} positive solutions. Let \(\varphi\) be any
positive solution of \eqref{eq:Lichnerowicz-exact-corrected}. Let
\[
m:=\min_{M}\varphi,
\qquad
M:=\max_{M}\varphi.
\]
At a point of maximum,
\[
0
=
-8\Delta_{g}\varphi + R[g]\varphi + \frac23\varphi^{5}-S\varphi^{-7}
\ge
R[g]M+\frac23 M^{5}-S M^{-7},
\]
since $\Delta_{g} \varphi<0$ at that point.
Hence
\[
\frac23 M(M^{4}-1)\le |\rho|\,M + S M^{-7}.
\]
Similarly, at a point of minimum,
\[
0
=
-8\Delta_{g}\varphi + R[g]\varphi + \frac23\varphi^{5}-S\varphi^{-7}
\le
R[g]m+\frac23 m^{5}-S m^{-7},
\]
since $\Delta_{g} \varphi>0$ at that point.
Hence
\[
\frac23 m(1-m^{4})\le |\rho|\,m + S m^{-7}.
\]
Using \eqref{eq:defect-small-proof}, one now chooses \(\delta_{0}\) so small that these two
inequalities force
\[
\frac34\le m\le M\le \frac54.
\]
Thus every positive solution takes values in \([\frac34,\frac54]\).

\noindent  Now let \(\varphi_{1},\varphi_{2}\) be two positive solutions, and set
\[
w:=\varphi_{1}-\varphi_{2}.
\]
Subtracting the two equations and using the mean value theorem pointwise yields
\[
-8\Delta_{g}w + a(x)w=0,
\]
where
\[
a(x)
=
R[g]
+
\frac{10}{3}\theta_{1}(x)^{4}
+
7S(x)\theta_{2}(x)^{-8},
\]
for some \(\theta_{1}(x),\theta_{2}(x)\in[\frac34,\frac54]\). Hence
\[
a(x)\ge -\|\rho\|_{C^{0}(g)}-\frac23+\frac{10}{3}\Bigl(\frac34\Bigr)^{4}.
\]
After possibly shrinking \(\delta_{0}\) once more, the right-hand side is bounded below by
a positive constant \(c_{*}>0\). Therefore
\[
-8\Delta_{g}w+c_{*}w\le 0
\qquad\text{and}\qquad
-8\Delta_{g}(-w)+c_{*}(-w)\le 0,
\]
since the order of $\varphi_{1}$ and $\varphi_{2}$ here does not matter.
So the maximum principle implies \(w\equiv0\). Thus the positive solution is unique.

\noindent Define
\[
g_{0}:=\varphi^{4}g,\qquad \Sigma_{0}:=\varphi^{-2}\Sigma.
\]
The standard conformal covariance identities in dimension \(3\) give
\[
R[g_{0}]
=
\varphi^{-5}\bigl(-8\Delta_{g}\varphi+R[g]\varphi\bigr),
\qquad
|\Sigma_{0}|_{g_{0}}^{2}
=
\varphi^{-12}|\Sigma|_{g}^{2}.
\]
Hence \eqref{eq:Lichnerowicz-exact-corrected} is equivalent to
\[
R[g_{0}]+\frac23=|\Sigma_{0}|_{g_{0}}^{2}.
\]
Since \(\Sigma\) is \(g\)-trace-free and \(g\)-divergence-free, the standard conformal
transformation law for TT tensors implies
\[
\tr_{g_{0}}\Sigma_{0}=0,
\qquad
\div_{g_{0}}\Sigma_{0}=0.
\]

\noindent It remains to prove \eqref{eq:T-change-corrected}. In dimension \(3\),
\[
\mathfrak T[\varphi^{4}g]
=
\mathfrak T[g]
-
2\varphi^{-1}\Bigl(\nabla^{2}\varphi-\frac13(\Delta_{g}\varphi)g\Bigr)
+
6\varphi^{-2}\Bigl(d\varphi\otimes d\varphi-\frac13|\nabla\varphi|_{g}^{2}g\Bigr).
\]
Therefore
\begin{align*}
\mathfrak T[g_{0}]-\mathfrak T[g]
&=
-2\varphi^{-1}\Bigl(\nabla^{2}\varphi-\frac13(\Delta_{g}\varphi)g\Bigr)\\
&\qquad
+
6\varphi^{-2}\Bigl(d\varphi\otimes d\varphi-\frac13|\nabla\varphi|_{g}^{2}g\Bigr).
\end{align*}
Since \(\varphi\in[\frac34,\frac54]\) pointwise and
\(\|\varphi-1\|_{H^{4}(g)}\le C\|\mathcal E(g,\Sigma)\|_{H^{2}(g)}\), uniform composition
estimates on the bounded-geometry class imply
\[
\|\varphi^{-1}\|_{H^{4}(g)}+\|\varphi^{-2}\|_{H^{4}(g)}\le C.
\]
Using that \(H^{2}(M)\) is an algebra in dimension \(3\),
\[
\left\|
\varphi^{-1}\Bigl(\nabla^{2}\varphi-\frac13(\Delta_{g}\varphi)g\Bigr)
\right\|_{H^{2}(g)}
\le
C\|\varphi-1\|_{H^{4}(g)},
\]
and
\[
\left\|
\varphi^{-2}\Bigl(d\varphi\otimes d\varphi-\frac13|\nabla\varphi|_{g}^{2}g\Bigr)
\right\|_{H^{2}(g)}
\le
C\|\varphi-1\|_{H^{4}(g)}^{2}.
\]
Hence
\[
\|\mathfrak T[g_{0}]-\mathfrak T[g]\|_{H^{2}(g)}
\le
C\|\varphi-1\|_{H^{4}(g)}
+
C\|\varphi-1\|_{H^{4}(g)}^{2}.
\]
For \(\delta_{0}\) sufficiently small, the quadratic term is absorbed into the linear one,
and \eqref{eq:phi-close-corrected} yields
\[
\|\mathfrak T[g_{0}]-\mathfrak T[g]\|_{H^{2}(g)}
\le
C\|\mathcal E(g,\Sigma)\|_{H^{2}(g)}.
\]
This proves \eqref{eq:T-change-corrected}.

\noindent Finally, since \(g\) and \(\Sigma\) are smooth and \(\varphi\) is bounded above and below
away from zero, standard elliptic bootstrapping applied to
\eqref{eq:Lichnerowicz-exact-corrected} shows that \(\varphi\in C^{\infty}(M)\).
\end{proof}

\end{document}
